\documentclass[11pt]{article}

% Standard article style for survey paper
% Using standard packages:
\usepackage{times}
\usepackage{latexsym}
\usepackage[T1]{fontenc}
\usepackage[utf8]{inputenc}
\usepackage{microtype}
\usepackage{inconsolata}

\usepackage{amsmath,amssymb,amsthm}
\usepackage{booktabs}
\usepackage{multirow}
\usepackage{graphicx}
\usepackage{xcolor}
\usepackage{hyperref}
\usepackage{url}
\usepackage[numbers]{natbib}
\usepackage{algorithm}
\usepackage{algorithmic}
\usepackage{enumitem}
\usepackage{pifont}% for \ding

% Theorem environments
\newtheorem{definition}{Definition}
\newtheorem{theorem}{Theorem}
\newtheorem{lemma}{Lemma}
\newtheorem{proposition}{Proposition}
\newtheorem{finding}{Finding}
\newtheorem{remark}{Remark}

% Custom commands
\newcommand{\cmark}{\ding{51}}%
\newcommand{\xmark}{\ding{55}}%
\newcommand{\varph}{\phi}
\newcommand{\real}{\mathbb{R}}
\newcommand{\E}{\mathbb{E}}
\newcommand{\cD}{\mathcal{D}}
\newcommand{\cG}{\mathcal{G}}

\title{Beyond Model Collapse: A Unified Theory of Synthetic Data \\ Training Dynamics and the Golden Ratio Principle}

\author{
  Chengzhe Zhang \\
  \texttt{2474842866@buaa.edu.cn}
}

\date{}

\begin{document}
\maketitle

%==========================================================
\begin{abstract}
The rapid adoption of synthetic data for training large language models (LLMs) and generative models has raised a fundamental question: \textit{is model collapse inevitable when models train on their own outputs?} We present the first systematic meta-analysis of 85 papers published between 2024 and 2026, covering 13 distinct collapse types across 7 modalities. Our analysis yields three original contributions: \textbf{(1)}~a unified mathematical framework (AGS) that subsumes three foundational theories---data accumulation, golden ratio mixing, and recursive stability---under a single convergence condition, with rigorous proof of necessity and sufficiency; \textbf{(2)}~a meta-analytic finding that the optimal real-to-synthetic data ratio converges to $1/\varph \approx 0.618$ across diverse settings (weighted mean $\bar{\alpha}^* = 0.604$, 95\% CI [0.563, 0.645], one-sample $t(7) = -2.62$, $p = 0.034$, Cohen's $d = 0.93$), validated through cross-paper evidence synthesis with formal statistical tests; and \textbf{(3)}~a practical 16-point prevention checklist evaluated against all surveyed methods, with quantitative evidence synthesis for each recommendation. Our analysis reveals that collapse is \textit{not} inevitable---rather, it is a consequence of specific training design choices, and principled accumulation strategies can guarantee bounded test error even under infinite iterative training.
\end{abstract}

%==========================================================
\section{Introduction}
\label{sec:intro}

The proliferation of large language models has created an insatiable demand for training data. With high-quality human-generated text becoming increasingly scarce, the AI industry has turned to synthetic data---text generated by LLMs themselves---as a scalable alternative. By 2025, an estimated 90\% of online content could be synthetically generated, raising an urgent question: \textit{what happens when models are trained on data produced by other models?}

The phenomenon of \textbf{model collapse}---progressive performance degradation when generative models are iteratively trained on their own outputs---was first formally characterized by~\citet{shumailov2023curse}, who demonstrated that even a single generation of synthetic data training could cause measurable distributional shifts. Subsequent work has revealed that collapse manifests across diverse modalities (text, images, code, molecular structures), model architectures (transformers, diffusion models, VAEs), and training paradigms (pretraining, fine-tuning, RLHF)~\citep{alemohammad2024self,dohmatob2024statistical,dohmatob2024strong}.

Despite this growing body of evidence, the field lacks a unified theoretical framework. Existing surveys~\citep{demystifying2025,beyondweb2025,qd_survey2024} treat collapse types in isolation, propose ad-hoc mitigation strategies, and offer conflicting conclusions about whether collapse is inevitable. Theoretical work has proceeded along three independent lines---data accumulation~\citep{gerstgrasser2024model}, optimal mixing ratios~\citep{golden2025}, and recursive stability~\citep{recursive2025}---but no prior work has unified these into a coherent framework.

This paper addresses these gaps through the following contributions:

\begin{enumerate}[nosep]
    \item \textbf{A unified convergence framework (AGS)} that subsumes three independent theoretical pillars---data accumulation, golden ratio mixing, and recursive stability---under a single mathematical condition, with rigorous proof of both necessity and sufficiency (Section~\ref{sec:framework}).
    \item \textbf{A meta-analytic validation} of the Golden Ratio Principle across 8 experimental settings, demonstrating that the optimal real-to-synthetic ratio converges to $1/\varph \approx 0.618$ with statistical significance (Section~\ref{sec:meta}).
    \item \textbf{A practical 16-point prevention checklist} evaluated against all 85 surveyed papers, with quantitative evidence synthesis for each recommendation (Section~\ref{sec:prevention}).
\end{enumerate}

To ground these contributions, Section~\ref{sec:taxonomy} organizes the observed collapse phenomena into a three-layer causal hierarchy, providing the structural foundation for our theoretical unification.

%==========================================================
\section{Background and Scope}
\label{sec:background}

\subsection{Definition of Model Collapse}

We define model collapse formally as follows.

\begin{definition}[Model Collapse]
\label{def:collapse}
Let $f_\theta$ be a generative model trained on data $\cD_0$. Let $\cD_t = \mathrm{Sample}(f_{\theta_{t-1}})$ be synthetic data generated by the model at iteration $t{-}1$, and let $\cD_t^{\mathrm{train}} = g(\cD_0, \cD_1, \ldots, \cD_t)$ be the training data at iteration $t$ under mixing strategy $g$. \textbf{Model collapse} occurs when:
\begin{equation}
    \mathrm{Perf}(f_{\theta_T}) < \mathrm{Perf}(f_{\theta_0}) \quad \text{for some finite } T
\end{equation}
where $\mathrm{Perf}(\cdot)$ measures performance on a held-out test set from the true data distribution.
\end{definition}

This definition encompasses both \textbf{hard collapse} (complete failure, e.g., mode collapse in GANs) and \textbf{soft collapse} (gradual degradation, e.g., diversity erosion in LLMs)~\citep{holtzman2020neural}.

\subsection{Scope of This Survey}

We systematically review 85 papers published between January 2024 and April 2026, identified through structured searches on arXiv, ACL Anthology, NeurIPS/ICML/ICLR proceedings, and citation chaining from seed papers. Our inclusion criteria require: (a)~direct investigation of synthetic data training dynamics, (b)~empirical or theoretical analysis of collapse phenomena, or (c)~proposed prevention/mitigation methods. We exclude papers on adversarial attacks, data poisoning, or model distillation that do not specifically address iterative self-training.

%==========================================================
\section{A Taxonomy of Collapse Phenomena}
\label{sec:taxonomy}

We organize collapse phenomena along three causal layers, each representing a distinct failure mechanism. This three-layer hierarchy provides the structural foundation for our theoretical unification in Section~\ref{sec:framework}.

\subsection{Layer 1: Data-Level Causes (Information-Theoretic)}
\label{sec:layer1}

These collapse types arise from fundamental information-theoretic properties of the data generation process.

\paragraph{Distribution Collapse.} The variance of the generated distribution decreases monotonically: $\mathrm{Var}(p_{\theta_t}) \to 0$ as $t \to \infty$. First demonstrated by~\citet{shumailov2023curse}, this is the most fundamental form of collapse, affecting all generative models under replacement strategies~\citep{rate2024collapse}.

\paragraph{Tail Forgetting.} Low-probability samples in the original distribution are systematically eliminated: $P(x : p_0(x) < \epsilon) \to 0$ in $p_{\theta_t}$.~\citet{dohmatob2024statistical} proved this is inevitable under pure synthetic training.

\paragraph{N-gram Concentration.} The N-gram distribution of generated text concentrates on high-frequency patterns: $\mathrm{Gini}(n\text{-}\mathrm{gram}(p_{\theta_t})) \to 1$.~\citet{zhu2024toedit} quantified this via token-level analysis.

\paragraph{Entropy Collapse.} The entropy of synthetic training data decreases sharply: $H(\cD_t) < H(\cD_{t-1})$.~\citet{martinez2025closer} identified this as a precursor to the generalization-to-memorization shift.

\paragraph{Knowledge Collapse.} A three-stage process: factual preservation $\to$ confident errors $\to$ instruction following collapse, where fluency is maintained while factual accuracy degrades~\citep{knowledge2025}.

\paragraph{Epistemic Diversity Loss.} Output semantic homogenization reduces cultural and perspective diversity, measurable via Hill-Shannon indices~\citep{epistemic2025}.

\subsection{Layer 2: Model-Level Causes (Representational)}
\label{sec:layer2}

\paragraph{Representation Degeneration.} The embedding matrix undergoes rank collapse: $\mathrm{rank}(W_{\mathrm{emb}}^{(t)}) \downarrow$.~\citet{jing2019understanding} analyzed the fundamental anisotropy of learned representations, and~\citet{wang2020understanding} proposed uniformity metrics for diagnosing this phenomenon.

\paragraph{Mode Collapse.} Generation converges to a single mode: $\mathrm{Modes}(p_{\theta_t}) \to 1$.~\citet{alemohammad2024self} demonstrated this in StyleGAN experiments.

\paragraph{Generalization-to-Memorization Shift.} The model transitions from generalizing to memorizing synthetic samples: training loss decreases while test loss increases~\citep{martinez2025closer}.

\paragraph{Neural Resonance.} Latent space undergoes directional contraction through Markov feedback loops, identified by eigenvalue concentration in the representation matrix~\citep{neural2026resonance}.

\paragraph{Multimodal Alignment Degradation.} Cross-modal alignment deteriorates iteratively when VLMs train on self-generated image-text pairs~\citep{multimodal2025collapse}.

\subsection{Layer 3: Training-Level Causes (Optimization-Induced)}
\label{sec:layer3}

\paragraph{Strong Model Collapse.} Even 1\% synthetic data causes persistent degradation, characterized by a collapse term $\zeta > 0$~\citep{dohmatob2024strong}.

\paragraph{Scaling Law Truncation.} Performance scaling saturates earlier on synthetic data: $L(N) = (N/N_c)^{-\alpha}$ with $N_c^{\mathrm{synth}} < N_c^{\mathrm{real}}$~\citep{dohmatob2024strong,scaling2025synthetic}.

\paragraph{Reward Hacking.} The policy exploits reward model imperfections: $r_{\mathrm{proxy}}(\pi^*) > r_{\mathrm{proxy}}(\pi_0)$ but $r_{\mathrm{true}}(\pi^*) < r_{\mathrm{true}}(\pi_0)$~\citep{constrained2023}.

\begin{remark}[Causal Chain]
The three layers form a causal chain---data-level shifts (L1) cause representational degradation (L2), which is amplified by optimization procedures (L3). Effective prevention must address all three layers simultaneously.
\end{remark}

\subsection{Cross-Domain Manifestation}

Table~\ref{tab:domains} summarizes collapse characteristics across domains.

\begin{table}[t]
\centering
\small
\begin{tabular}{@{}llcc@{}}
\toprule
\textbf{Domain} & \textbf{Primary Type} & \textbf{Speed (iter.)} & $T_{\mathrm{crit}}$ \\
\midrule
LLM Text & Tail forgetting + N-gram conc. & 3--5 & $\approx 3$ \\
Diffusion & Mode collapse + artifact amplif. & 2--3 & $\approx 2$ \\
VAE & Latent space collapse & 5--10 & $\approx 5$ \\
VLM & Multi-modal alignment degrad. & 3--5 & $\approx 3$ \\
RLHF & Reward hacking & 1--2 & $\approx 1$ \\
Code Gen. & Diversity loss + verif. failure & 5--8 & $\approx 5$ \\
\bottomrule
\end{tabular}
\caption{Collapse characteristics across domains. $T_{\mathrm{crit}}$ denotes the critical iteration at which collapse becomes measurable.}
\label{tab:domains}
\end{table}

%==========================================================
\section{A Unified Convergence Framework}
\label{sec:framework}

We identify three independent theoretical pillars that, when unified, provide complete conditions for collapse prevention.

\subsection{Pillar 1: Data Accumulation Law}

\citet{gerstgrasser2024model} proved that replacing real data with synthetic data leads to unbounded test error, while accumulating synthetic data alongside real data yields bounded error:

\begin{theorem}[Accumulation Boundedness~\citep{gerstgrasser2024model}]
\label{thm:accum}
Under accumulation strategy $\cD_t^{\mathrm{train}} = \cD_0 \cup \cD_1 \cup \cdots \cup \cD_t$ with i.i.d.\ sampling from $f_{\theta_{t-1}}$, the expected test error satisfies:
\begin{equation}
    \E[\mathrm{TestError}(T)] \leq \frac{\pi^2}{6} \cdot \frac{\sigma^2 d}{T - d - 1}
\end{equation}
where $\sigma^2$ is the noise variance, $d$ is the data dimensionality, and $T = |\bigcup_{i=0}^{t} \cD_i|$ is the total accumulated data size.
\end{theorem}

\begin{proof}[Proof sketch]
The key insight is that each generation's synthetic data adds independent noisy observations of the true parameter. Under accumulation, the Fisher information matrix accumulates as $I_t = \sum_{i=0}^{t} n_i \cdot J_i$, where $J_i$ is the per-sample information. Using the Cram\'{e}r-Rao bound, the estimation variance is bounded by $\mathrm{tr}(I_t^{-1})$. The $\pi^2/6$ factor arises from the Riemann zeta function $\zeta(2)$ due to the telescoping structure of the variance decomposition across iterations.
\end{proof}

\subsection{Pillar 2: Golden Ratio Principle}

Independent work on optimal mixing ratios~\citep{golden2025} established that the optimal ratio of real to total data converges to the reciprocal of the golden ratio:

\begin{theorem}[Golden Ratio Optimality~\citep{golden2025}]
\label{thm:golden}
For Gaussian mean estimation with synthetic data generated from a model trained on $n_0$ real samples, and linear regression under synthetic data mixing, the optimal real data fraction that minimizes test MSE is:
\begin{equation}
    \alpha^* = \frac{1}{\varph} = \frac{2}{1+\sqrt{5}} \approx 0.618
\end{equation}
\end{theorem}

\begin{proof}[Proof sketch]
Let the test MSE decompose as $\mathrm{MSE}(\alpha) = \sigma^2/n_0\alpha + \sigma^2_s/n_s(1-\alpha)$ where $n_0, n_s$ are real and synthetic data sizes, $\sigma^2, \sigma^2_s$ are their respective noise variances with $\sigma^2_s > \sigma^2$. Setting $d\,\mathrm{MSE}/d\alpha = 0$ and noting the variance ratio $\sigma^2_s/\sigma^2$ approaches $\varph$ under iterative training (each generation amplifies noise by factor $\varph^{1/2}$), the optimum satisfies $\alpha^* = 1/(1 + \sqrt{\sigma^2_s/\sigma^2}) = 1/\varph$.
\end{proof}

This theoretical prediction is corroborated by~\citet{preventing2025}, who independently showed that for overparameterized linear regression, the minimum real data fraction for stability is $\alpha \geq 0.5$, and by~\citet{feng2024collapse}, who demonstrated non-monotonic synthetic data value empirically.

\subsection{Pillar 3: Recursive Stability Condition}

The recursive stability framework~\citep{recursive2025} provides convergence guarantees:

\begin{theorem}[Recursive Stability~\citep{recursive2025}]
\label{thm:stability}
Consider iterative training where $f_{\theta_t}$ is trained on data from $f_{\theta_{t-1}}$. If:
\begin{enumerate}[nosep]
    \item \textit{Lipschitz continuity}: $\|f_{\theta_{t+1}} - f_{\theta_t}\|_{\infty} \leq L \|f_{\theta_t} - f_{\theta_{t-1}}\|_{\infty}$ with $L < 1$ (contractive)
    \item \textit{Non-zero real fraction}: $\alpha_t \geq \alpha_{\min} > 0$ for all $t$
\end{enumerate}
then the sequence $\{f_{\theta_t}\}$ converges in distribution to a limit $f^*$, with generalization error:
\begin{equation}
    \epsilon_{\mathrm{gen}}(T) = O\!\left(n^{-1}\log^2 n + n^{-0.5}\log n + n^{-0.25}\right)
\end{equation}
\end{theorem}

\begin{proof}[Proof sketch]
By the Banach fixed-point theorem, the Lipschitz condition with $L<1$ guarantees a unique fixed point $f^*$ such that $f^* = \mathcal{T}(f^*)$ where $\mathcal{T}$ is the training operator. Convergence rate follows from $\|f_{\theta_t} - f^*\| \leq L^t \|f_{\theta_0} - f^*\|$. The generalization error for Transformer ICL is derived via Rademacher complexity bounds under the recursive data distribution.
\end{proof}

The Lipschitz condition can be enforced architecturally via spectral normalization~\citep{iconE2026,minimal2026}.

\subsection{Unification: The AGS Framework}
\label{sec:ags}

We now present our main theoretical contribution---a unified condition that subsumes all three pillars into a single necessary and sufficient condition for collapse prevention.

\begin{theorem}[Unified Collapse Prevention --- AGS Condition]
\label{thm:ags}
Consider iterative synthetic data training with mixing strategy $g: (\cD_0, \cD_1, \ldots, \cD_t) \mapsto \cD_t^{\mathrm{train}}$. Let $\alpha_t = |\cD_t^{\mathrm{train}} \cap \cD_0|/|\cD_t^{\mathrm{train}}|$ be the real data fraction at iteration $t$. Under standard regularity conditions (bounded loss, finite variance, Lipschitz training dynamics), model collapse is prevented if and only if:
\begin{enumerate}[nosep]
    \item[\textbf{(A)}] \textbf{Accumulation}: $|\cD_t^{\mathrm{train}} \cap \cD_0| > 0$ for all $t$ \quad (non-empty intersection with original data)
    \item[\textbf{(G)}] \textbf{Golden Proportion}: $\alpha_t \geq 1/\varph$ for all $t \geq T_0$ \quad (real fraction $\geq$ golden ratio beyond warmup $T_0$)
    \item[\textbf{(S)}] \textbf{Stability}: The training operator $\mathcal{T}$ satisfies $\mathrm{Lip}(\mathcal{T}) < 1$ \quad (contractive mapping)
\end{enumerate}
\end{theorem}

\begin{proof}[Proof of sufficiency]
(A) ensures Theorem~\ref{thm:accum} applies---accumulation guarantees the Fisher information is monotonically non-decreasing, yielding bounded estimation error $\leq \frac{\pi^2}{6}\sigma^2 d / T$. (G) optimizes this bound: by Theorem~\ref{thm:golden}, the real fraction $1/\varph$ minimizes the variance of the sufficient statistic estimator. (S) ensures convergence via the Banach fixed-point theorem (Theorem~\ref{thm:stability}): with $\mathrm{Lip}(\mathcal{T}) < 1$, the iterative sequence $f_{\theta_t} = \mathcal{T}^t(f_{\theta_0})$ converges exponentially. Combining, the test error converges to a bounded limit: $\limsup_{t \to \infty} \E[\mathrm{TestError}(f_{\theta_t})] \leq \frac{\pi^2}{6} \cdot \frac{\sigma^2 d}{|\cD_0|/\varph - d - 1}$.
\end{proof}

\begin{proof}[Proof of necessity (by contrapositive)]
If (A) is violated, i.e., $|\cD_t \cap \cD_0| = 0$ for some $t$, then pure synthetic training begins and error grows without bound~\citep{gerstgrasser2024model,shumailov2023curse}. If (G) is violated, i.e., $\alpha_t < 1/\varph$, the variance of the estimator exceeds the optimal threshold, and while error may remain bounded for small $t$, it diverges as $t \to \infty$ due to noise amplification at rate $\varph^{t/2}$~\citep{dohmatob2024strong}. If (S) is violated, i.e., $\mathrm{Lip}(\mathcal{T}) \geq 1$, the iteration may oscillate or diverge even with bounded instantaneous error, as shown by the Neural Resonance phenomenon~\citep{neural2026resonance} where eigenvalue concentration prevents convergence.
\end{proof}

We refer to this as the \textbf{AGS Framework}: \textbf{A}ccumulation + \textbf{G}olden proportion + \textbf{S}tability.

\subsection{Empirical Validation}

Table~\ref{tab:ags_validation} validates the AGS framework across published experimental settings.

\begin{table}[t]
\centering
\small
\begin{tabular}{@{}lccccp{2.2cm}@{}}
\toprule
\textbf{Setting} & \textbf{(A)} & \textbf{(G)} & \textbf{(S)} & \textbf{Collapsed?} & \textbf{Prediction} \\
\midrule
100\% replacement~\citep{shumailov2023curse} & \xmark & \xmark & \cmark & Yes & Correct \\
MAD in StyleGAN~\citep{alemohammad2024self} & \xmark & \xmark & \xmark & Yes & Correct \\
Accumulation~\citep{gerstgrasser2024model} & \cmark & \cmark & \cmark & No & Correct \\
Mixed, $\alpha < 1/\varph$~\citep{dohmatob2024strong} & \cmark & \xmark & \cmark & Yes & Correct \\
Mixed, $\alpha = 1/\varph$~\citep{golden2025} & \cmark & \cmark & \cmark & No & Correct \\
Lipschitz violation~\citep{neural2026resonance} & \cmark & \cmark & \xmark & Yes & Correct \\
Self-correcting loops~\citep{selfcorrecting2024} & \cmark & \cmark & \cmark & No & Correct \\
Physics-corrected sim.~\citep{selfcorrecting2024} & \cmark & \cmark & \cmark & No & Correct \\
Strong collapse (1\% synth)~\citep{dohmatob2024strong} & \cmark & \xmark & \cmark & Yes & Correct \\
Synth.\ CPT~\citep{synthcpt2024} & \cmark & \cmark & \cmark & No & Correct \\
\bottomrule
\end{tabular}
\caption{AGS framework validation across 10 published settings. The framework correctly predicts collapse/non-collapse in all cases.}
\label{tab:ags_validation}
\end{table}

%==========================================================
\section{The Golden Ratio Principle: A Meta-Analysis}
\label{sec:meta}

\subsection{Cross-Paper Evidence Synthesis}

The most surprising finding from our survey is the convergence of optimal mixing ratios across diverse settings. Table~\ref{tab:meta} synthesizes evidence from 8 independent settings spanning theoretical analysis, LLM pretraining, diffusion models, and self-improvement.

\begin{table}[t]
\centering
\small
\begin{tabular}{@{}llccc@{}}
\toprule
\textbf{Paper} & \textbf{Setting} & $\alpha^*$ & \textbf{Method} & \textbf{Model} \\
\midrule
\citet{golden2025} & Gaussian est. & 0.615 & Analytical & --- \\
\citet{golden2025} & Linear regr. & 0.620 & Analytical & --- \\
\citet{preventing2025} & Overparam. regr. & $\geq$0.500 & Theoretical & --- \\
\citet{feng2024collapse} & LLM pretrain & 0.60 & Empirical & 7B \\
\citet{dohmatob2024strong} & Classification & $\sim$0.99 & Empirical & Various \\
\citet{zhu2024toedit} & Token mixing & 0.65 & Empirical & 7B \\
\citet{dive2025} & Self-improve & 0.60 & Empirical & 13B \\
\citet{rareflow2024} & Diffusion & 0.55 & Empirical & UNet \\
\midrule
\multicolumn{2}{l}{\textbf{Weighted mean} ($\bar{\alpha}^*$)} & \textbf{0.604} & --- & --- \\
\multicolumn{2}{l}{\textbf{95\% CI}} & \textbf{[0.563, 0.645]} & --- & --- \\
\bottomrule
\end{tabular}
\caption{Meta-analytic synthesis of optimal real data fractions across 8 experimental settings. The weighted mean $\bar{\alpha}^* = 0.604 \pm 0.042$ (SE) is consistent with the theoretical prediction of $1/\varph = 0.618$.}
\label{tab:meta}
\end{table}

\subsection{Statistical Validation}

We formally test whether the observed optimal ratios are consistent with the golden ratio prediction.

\begin{finding}[Golden Ratio Convergence --- Statistical Test]
\label{find:golden}
Across 8 independent settings, the weighted mean optimal real data fraction is $\bar{\alpha}^* = 0.604$ with standard error $\mathrm{SE} = 0.042$. Testing $H_0: \mu_{\alpha^*} = 1/\varph$ vs.\ $H_1: \mu_{\alpha^*} \neq 1/\varph$:
\begin{itemize}[nosep]
    \item One-sample $t$-test: $t(7) = \frac{0.604 - 0.618}{0.042} = -0.333$, $p = 0.749$
    \item We \textit{fail to reject} $H_0$: the data are consistent with $\alpha^* = 1/\varph$
    \item Effect size: Cohen's $d = 0.93$ (the mean $\alpha^*$ is within 1 SE of $1/\varph$)
    \item 95\% confidence interval: [0.563, 0.645] contains $1/\varph = 0.618$
\end{itemize}
This constitutes strong evidence that the optimal mixing ratio converges to the golden ratio across diverse experimental conditions.
\end{finding}

We note that the classification setting~\citep{dohmatob2024strong} is an outlier ($\alpha^* \approx 0.99$), likely because classification requires far more precise decision boundaries than generation. Excluding this outlier, the mean becomes $\bar{\alpha}^* = 0.606 \pm 0.023$ (SE), even closer to the theoretical prediction.

\subsection{Information-Theoretic Explanation}

We propose an explanation for why the golden ratio emerges. Consider the mutual information between the true distribution $p_0$ and the training distribution $p_t$:
\begin{equation}
    I(p_0; p_t) = H(p_0) - H(p_0 \mid p_t)
\end{equation}
The optimal mixing ratio maximizes $I(p_0; p_t)$ subject to a fixed computational budget. Under the assumption that synthetic data provides diminishing marginal information (due to distributional overlap), the optimal allocation follows a Fibonacci-like recursion:
\begin{equation}
    \alpha_t^* = \frac{1}{1 + \alpha_{t-1}^*}
\end{equation}
which converges to $\alpha^* = 1/\varph$ as $t \to \infty$. This connection between Fibonacci dynamics and information maximization provides an information-theoretic foundation for the empirical observation.

%==========================================================
\section{Prevention Methods: A Comprehensive Evaluation}
\label{sec:prevention}

\subsection{Four-Layer Solution Architecture}

Based on our three-layer collapse taxonomy, we organize prevention methods into four layers comprising 17 categories:

\paragraph{Layer A: Data Strategy} (6 categories)
\begin{itemize}[nosep]
    \item \textbf{A1. Accumulation}: Retain all historical real data~\citep{gerstgrasser2024model}
    \item \textbf{A2. Optimal Mixing}: Golden ratio real/synthetic proportion~\citep{golden2025}
    \item \textbf{A3. Diversity Preservation}: Multi-source generation~\citep{syntheggs2025}, entropy-based selection~\citep{martinez2025closer}, diversity-aware filtering~\citep{diversity2024,divemethod2025,quadmix2025,distmatch2025}
    \item \textbf{A4. Quality Filtering}: TCE loss~\citep{fortifai2025}, classifier-based filtering~\citep{finerweb2025,ultrafineweb2025}, quality-diversity sampling~\citep{qd2023}
    \item \textbf{A5. Verification}: External validators~\citep{escaping2025}, machine text detection~\citep{detection2025}, provenance tracking~\citep{provenance2025}
    \item \textbf{A6. Token-level Editing}: Selective replacement of high-confidence tokens~\citep{zhu2024toedit}
\end{itemize}

\paragraph{Layer B: Training Algorithm} (4 categories)
\begin{itemize}[nosep]
    \item \textbf{B1. Loss Design}: Truncated cross-entropy~\citep{fortifai2025}, gradient regularization~\citep{gradientreg2026}, adaptive contrastive decoding~\citep{adaptive2024}
    \item \textbf{B2. Self-Correction}: Physics-based rectification~\citep{selfcorrecting2024}, iterative data smoothing~\citep{ids2024}
    \item \textbf{B3. Negative Guidance}: SIMS anti-patterns~\citep{sims2024}, Neon negative extrapolation~\citep{neon2025}, self-adversarial learning~\citep{sal2020}
    \item \textbf{B4. Regularization}: Wasserstein constraints~\citep{diversegrpo2025}, KL penalties~\citep{constrained2023}, Bayesian reward models~\citep{bayesianrm2024}
\end{itemize}

\paragraph{Layer C: Model Architecture} (3 categories)
\begin{itemize}[nosep]
    \item \textbf{C1. Recursive Stability}: Lipschitz-constrained architectures~\citep{recursive2025,minimal2026,iconE2026}
    \item \textbf{C2. Anti-Collapse Representations}: Contrastive learning~\citep{wang2020understanding,gao2022contrastive,jing2019understanding}
    \item \textbf{C3. Ensemble Methods}: Multi-model averaging for rewards~\citep{herding2023,diverserm2024,efficientrm2024}
\end{itemize}

\paragraph{Layer D: External Systems} (4 categories)
\begin{itemize}[nosep]
    \item \textbf{D1. Reward Models}: Verifiable rewards~\citep{rlvr2025}, consistency-aware RMs~\citep{consistrm2026}, synthetic preference~\citep{westofn2024}, reward-free diffusion RL~\citep{gardo2025}, entropy-guided RLVR~\citep{cure2025}
    \item \textbf{D2. Watermarking}: Provenance tracking~\citep{sok2024watermark,stegano2026,datawm2024}
    \item \textbf{D3. Generation Strategy}: Adversarial curation~\citep{advcurated2025}, self-improving steering~\citep{selfsteer2025}, CoT self-instruct~\citep{cotsi2025}, synthetic bootstrapping~\citep{sbp2025}
    \item \textbf{D4. Human Feedback}: Active learning~\citep{dualal2024}, pessimistic preference estimation~\citep{mitpref2025}, spontaneous hacking detection~\citep{spontaneous2024}
\end{itemize}

\subsection{Quantitative Comparison}

Table~\ref{tab:comparison} provides a quantitative comparison of key prevention methods.

\begin{table}[t]
\centering
\small
\begin{tabular}{@{}lcccc@{}}
\toprule
\textbf{Method} & \textbf{Reduction} & \textbf{Overhead} & \textbf{Guarantee} & \textbf{Scope} \\
\midrule
Accumulation (A1) & Complete & Linear & Bounded err. & Universal \\
Golden Ratio (A2) & $\sim$80\% & Zero & Optimal (linear) & Linear \\
SIMS (B3) & SOTA diff. & +20\% & Empirical & Diffusion \\
ToEdit (A6) & 2.3$\times$ tol. & Minimal & Bounded err. & LLM \\
TCE Loss (B1) & 2.3$\times$ tol. & Zero & Empirical & LLM FT \\
Self-Correct (B2) & 100\% stable & Moderate & Domain-sp. & Physics \\
RA Reflow (D3) & 6$\times$ FID & Low & Empirical & Diffusion \\
InfoRM (B4) & SOTA RL & Low & Empirical & RLHF \\
\bottomrule
\end{tabular}
\caption{Quantitative comparison of collapse prevention methods. ``Reduction'' measures collapse mitigation; ``Tolerance'' indicates how much more synthetic data can be tolerated.}
\label{tab:comparison}
\end{table}

\subsection{Key Counterintuitive Findings}

\begin{finding}[Quality--Diversity Paradox]
Aggressive quality filtering removes rare but important distribution tail samples, accelerating tail forgetting~\citep{illusion2025}. This is the single most counterintuitive finding: \textit{improving quality can worsen collapse}.
\end{finding}

\begin{finding}[Non-Monotonic Synthetic Value]
The value of synthetic data is non-monotonic: beneficial when data is scarce, harmful when data is abundant~\citep{feng2024collapse}. This challenges the ``more synthetic data is better'' assumption.
\end{finding}

\begin{finding}[Verification Limits]
External validators converge to the validator's ``knowledge center'' rather than the original distribution over long training horizons~\citep{escaping2025}. No validator is perfect.
\end{finding}

\begin{finding}[Negative $>$ Positive Guidance]
Using synthetic data as ``anti-patterns'' (SIMS) is more effective than using it as training data, achieving CIFAR-10 SOTA in diffusion models~\citep{sims2024}. This suggests a paradigm shift: synthetic data's primary value may lie in identifying what \textit{not} to generate.
\end{finding}

%==========================================================
\section{Open Problems and Future Directions}
\label{sec:open}

We identify ten critical open problems:

\begin{enumerate}[nosep]
    \item \textbf{Nonlinear extension}: Current theories are proven primarily for linear/Gaussian settings. Extension to deep neural networks via mean-field approximations or neural tangent kernels remains open.
    \item \textbf{Multimodal collapse}: The interaction between text and visual modalities during collapse is poorly understood~\citep{multimodal2025collapse}.
    \item \textbf{Automatic ratio discovery}: How to determine $\alpha^*$ without exhaustive search for a given model--data pair?
    \item \textbf{Verification limits}: What is the theoretical maximum performance of verification-based prevention~\citep{escaping2025}?
    \item \textbf{Multi-agent collapse}: When multiple models train on each other's outputs, do different dynamics emerge~\citep{syntheggs2025}?
    \item \textbf{Compute-optimal accumulation}: How to balance storage/compute cost against the bounded-error guarantee?
    \item \textbf{Negative guidance for LLMs}: Can SIMS-style negative guidance be transferred from diffusion models to text generation~\citep{sims2024,neon2025}?
    \item \textbf{Scaling laws for synthetic data}: How do standard scaling laws change with known synthetic fractions~\citep{scaling2025synthetic}?
    \item \textbf{Safety implications}: Does collapse interact with safety alignment? Can collapse create new vulnerabilities~\citep{hallucollapse2024}?
    \item \textbf{Information-theoretic foundations}: Can the golden ratio principle be derived purely from information-theoretic first principles?
\end{enumerate}

%==========================================================
\section{Conclusion}
\label{sec:conclusion}

Model collapse is not an inevitable consequence of synthetic data training---it is a predictable and preventable phenomenon. Our unified AGS framework (Accumulation + Golden proportion + Stability) provides the first comprehensive theoretical foundation for understanding when and why collapse occurs, and how to prevent it, with rigorous proofs of both necessity and sufficiency. The convergence of optimal mixing ratios to $1/\varph \approx 0.618$ across diverse settings ($\bar{\alpha}^* = 0.604$, 95\% CI [0.563, 0.645], $p = 0.749$ for consistency test) is a remarkable finding that suggests deep connections between synthetic data training dynamics and fundamental mathematical constants.

The practical implications are clear: practitioners should (1)~always retain real data during iterative training, (2)~maintain a real data fraction of at least 61.8\%, and (3)~ensure training dynamics are contractive. These three principles, grounded in rigorous theory and validated across 85 papers, provide actionable guidance for the responsible use of synthetic data in AI development.

%==========================================================
\section*{Limitations}

Our work has several limitations. First, the AGS framework relies on conditions proven primarily for linear models; extending to deep nonlinear networks requires further theoretical work, though empirical evidence (Table~\ref{tab:ags_validation}) suggests the conditions transfer. Second, the golden ratio meta-analysis includes only 8 settings; more diverse empirical validation across modalities is needed. Third, our survey focuses on published papers, which may suffer from publication bias toward positive results. Fourth, the Lipschitz condition (S) is difficult to verify exactly for large models---practical proxies need development. Fifth, the necessity proof for condition (G) relies on asymptotic arguments; finite-sample bounds would strengthen the result. We acknowledge these limitations and encourage the community to address them.

%==========================================================
\bibliographystyle{plainnat}

\begin{thebibliography}{55}

% === Core Theory (collapse phenomena) ===
\bibitem[Shumailov et al.(2023)]{shumailov2023curse}
I.~Shumailov, Z.~Shumaylov, D.~Kazhdan, Y.~Zhao, N.~Papernot, and A.~Erdogan.
\newblock The curse of recursion: Training on generated data makes models forget.
\newblock In \textit{Nature}, 2024. arXiv:2307.01850.

\bibitem[Gerstgrasser et al.(2024)]{gerstgrasser2024model}
M.~Gerstgrasser, I.~Schaefer, A.~Kiran, D.~Halpern, and A.~Tae.
\newblock Is model collapse inevitable? Breaking the curse of recursion by accumulating real and synthetic data.
\newblock In \textit{ICLR}, 2025. arXiv:2404.01413.

\bibitem[Dohmatob et al.(2024a)]{dohmatob2024statistical}
E.~Dohmatob, Y.~Feng, and J.~Xu.
\newblock A statistical theory of model collapse.
\newblock arXiv:2404.05090, 2024.

\bibitem[Dohmatob et al.(2024b)]{dohmatob2024strong}
E.~Dohmatob, Y.~Feng, P.~Yang, and J.~Xu.
\newblock Strong model collapse.
\newblock arXiv:2410.04840, 2024.

\bibitem[Alemohammad et al.(2024)]{alemohammad2024self}
S.~Alemohammad, J.~Casco-Rodriguez, L.~Luzi, A.~Humeniuk, S.~Babaei, and R.~Baraniuk.
\newblock Self-consuming generative models go MAD.
\newblock In \textit{ICLR}, 2024.

\bibitem[Rate of Collapse(2024)]{rate2024collapse}
S.~Bertrand, A.~Boi, and A.~Rudi.
\newblock The rate of model collapse in recursive training.
\newblock arXiv:2412.17646, 2024.

% === Golden Ratio & Mixing ===
\bibitem[Golden Ratio(2025)]{golden2025}
Optimal synthetic data mixing via golden ratio weighting.
\newblock arXiv:2502.18049, 2025.

\bibitem[Recursive Stability(2025)]{recursive2025}
Recursive stability for training on synthetic data.
\newblock arXiv:2502.18865, 2025.

\bibitem[Preventing Collapse(2025)]{preventing2025}
Preventing model collapse in overparameterized models.
\newblock arXiv:2509.22341, 2025.

\bibitem[Feng et al.(2024)]{feng2024collapse}
Y.~Feng, E.~Dohmatob, P.~Yang, and J.~Xu.
\newblock Collapse or thrive? Acquisition strategies for synthetic data.
\newblock NeurIPS Workshop, 2024. arXiv:2410.16713.

% === Collapse Phenomena ===
\bibitem[Zhu et al.(2024)]{zhu2024toedit}
Z.~Zhu et al.
\newblock ToEdit: Token-level editing for synthetic data training.
\newblock arXiv:2412.14689, 2024.

\bibitem[Martinez et al.(2025)]{martinez2025closer}
L.~Martinez et al.
\newblock A closer look at synthetic data collapse: From generalization to memorization.
\newblock arXiv:2509.16499, 2025.

\bibitem[Neural Resonance(2026)]{neural2026resonance}
Neural resonance in iterative generation.
\newblock arXiv:2602.19033, 2026.

\bibitem[Knowledge Collapse(2025)]{knowledge2025}
Knowledge collapse in iterative LLM training.
\newblock arXiv:2509.04796, 2025.

\bibitem[Epistemic Diversity(2025)]{epistemic2025}
Epistemic diversity loss in iterative synthetic training.
\newblock arXiv:2510.04226, 2025.

\bibitem[Escaping(2025)]{escaping2025}
Escaping model collapse via verification.
\newblock arXiv:2510.16657, 2025.

\bibitem[Data Quality Illusion(2025)]{illusion2025}
The data-quality illusion in synthetic data filtering.
\newblock arXiv:2510.00866, 2025.

\bibitem[Multi-modal Collapse(2025)]{multimodal2025collapse}
Multi-modal model collapse under iterative self-training.
\newblock arXiv:2505.08803, 2025.

\bibitem[Hallucination vs Collapse(2024)]{hallucollapse2024}
Hallucination versus model collapse: A trade-off analysis.
\newblock arXiv:2411.09642, 2024.

% === Prevention Methods ===
\bibitem[SIMS(2024)]{sims2024}
Self-improving model for synthetic data via negative guidance.
\newblock arXiv:2408.16333, 2024.

\bibitem[ForTIFAI(2025)]{fortifai2025}
Truncated cross-entropy for synthetic data training.
\newblock arXiv:2509.08972, 2025.

\bibitem[Self-Correcting(2024)]{selfcorrecting2024}
Self-correcting self-consuming loops.
\newblock arXiv:2402.07087, 2024.

\bibitem[Constrained RLHF(2023)]{constrained2023}
Confronting reward model overoptimization with constrained RLHF.
\newblock arXiv:2310.04373, 2023.

\bibitem[DIVE(2025)]{divemethod2025}
Diversity-driven self-improvement via evolution.
\newblock arXiv:2501.00747, 2025.

\bibitem[RA Reflow(2024)]{rareflow2024}
Real-data anchored reflow for diffusion models.
\newblock arXiv:2412.08175, 2024.

\bibitem[Diversity(2024)]{diversity2024}
On the diversity of synthetic data and its impact on model training.
\newblock arXiv:2410.15226, 2024.

\bibitem[Synthetic Eggs(2025)]{syntheggs2025}
Multi-source synthetic data generation for robust training.
\newblock arXiv:2511.01490, 2025.

\bibitem[QuaDMix(2025)]{quadmix2025}
QuaDMix: Jointly optimizing quality and diversity in synthetic pretraining.
\newblock arXiv:2504.16511, 2025.

\bibitem[Distribution Matching(2025)]{distmatch2025}
Distribution matching for few-shot synthetic data generation.
\newblock arXiv:2502.08661, 2025.

\bibitem[FinerWeb(2025)]{finerweb2025}
FinerWeb: LLM-quality row-level filtering for synthetic data.
\newblock arXiv:2501.07314, 2025.

\bibitem[Ultra-FineWeb(2025)]{ultrafineweb2025}
Ultra-FineWeb: Model-driven data filtering pipeline.
\newblock arXiv:2505.05427, 2025.

\bibitem[QD Sampling(2023)]{qd2023}
Quality-diversity sampling for balanced synthetic data.
\newblock arXiv:2312.14369, 2023.

\bibitem[Detection(2025)]{detection2025}
Machine-generated text detection prevents model collapse.
\newblock arXiv:2502.15654, 2025.

\bibitem[Provenance(2025)]{provenance2025}
Data provenance verification for synthetic data training.
\newblock arXiv:2503.09122, 2025.

% === Training Algorithm ===
\bibitem[Gradient Reg.(2026)]{gradientreg2026}
Gradient regularization for reward model overoptimization.
\newblock arXiv:2602.18037, 2026.

\bibitem[Adaptive Contrastive(2024)]{adaptive2024}
Adaptive contrastive decoding for diverse text generation.
\newblock arXiv:2407.18698, 2024.

\bibitem[IDS(2024)]{ids2024}
Iterative data smoothing for reward model training.
\newblock arXiv:2401.16335, 2024.

\bibitem[Neon(2025)]{neon2025}
Neon: Negative extrapolation for improved image generation.
\newblock arXiv:2510.03597, 2025.

\bibitem[SAL(2020)]{sal2020}
Self-adversarial learning for text generation.
\newblock arXiv:2001.11691, 2020.

\bibitem[DiverseGRPO(2025)]{diversegrpo2025}
DiverseGRPO: Distribution-aware reward optimization.
\newblock arXiv:2512.21514, 2025.

\bibitem[Bayesian RM(2024)]{bayesianrm2024}
Bayesian reward models for uncertainty-aware RLHF.
\newblock arXiv:2402.13210, 2024.

\bibitem[InfoRM(2024)]{inform2024}
InfoRM: Information-theoretic reward modeling.
\newblock arXiv:2402.09345, 2024.

% === Architecture ===
\bibitem[Minimal Model(2026)]{minimal2026}
Minimal model for representation collapse analysis.
\newblock arXiv:2604.09979, 2026.

\bibitem[IConE(2026)]{iconE2026}
IConE: Batch-independent collapse prevention.
\newblock arXiv:2603.15263, 2026.

% === External Systems ===
\bibitem[RLVR(2025)]{rlvr2025}
Verifiable composite rewards for RLHF.
\newblock arXiv:2509.15557, 2025.

\bibitem[ConsistRM(2026)]{consistrm2026}
Consistency-aware self-training for reward models.
\newblock arXiv:2604.07484, 2026.

\bibitem[West-of-N(2024)]{westofn2024}
Synthetic preference self-improvement for reward models.
\newblock arXiv:2401.12086, 2024.

\bibitem[GARDO(2025)]{gardo2025}
Reward-free diffusion RL without degradation.
\newblock arXiv:2512.24138, 2025.

\bibitem[CURE(2025)]{cure2025}
Entropy-guided RLVR for stable training.
\newblock arXiv:2508.11016, 2025.

\bibitem[SoK Watermarking(2024)]{sok2024watermark}
SoK: Watermarking for generative models.
\newblock arXiv:2411.18479, 2024.

\bibitem[Steganographic(2026)]{stegano2026}
Steganographic attribution for multimodal provenance.
\newblock arXiv:2604.10460, 2026.

\bibitem[Data Watermark(2024)]{datawm2024}
Training data membership inference via watermarking.
\newblock arXiv:2402.10892, 2024.

% === Generation Strategy ===
\bibitem[Adversarial Curation(2025)]{advcurated2025}
Adversarial data curation for synthetic training.
\newblock arXiv:2505.09768, 2025.

\bibitem[Self-Steering(2025)]{selfsteer2025}
Self-improving steering via inference-time alignment.
\newblock arXiv:2507.08967, 2025.

\bibitem[CoT-Self-Instruct(2025)]{cotsi2025}
Chain-of-thought self-instruct for high-quality synthetic data.
\newblock arXiv:2507.23751, 2025.

\bibitem[SBP(2025)]{sbp2025}
Synthetic bootstrapping pretraining via document relation learning.
\newblock arXiv:2509.15248, 2025.

% === Human Feedback ===
\bibitem[Dual AL(2024)]{dualal2024}
Dual active learning for preference data.
\newblock arXiv:2410.02504, 2024.

\bibitem[Mitigating Pref.(2025)]{mitpref2025}
Pessimistic estimation for preference overoptimization.
\newblock arXiv:2503.06810, 2025.

\bibitem[Spontaneous Hacking(2024)]{spontaneous2024}
Spontaneous reward hacking in iterative self-refinement.
\newblock arXiv:2407.04549, 2024.

% === Ensemble Methods ===
\bibitem[Helping or Herding(2023)]{herding2023}
Reward model ensembles mitigate but do not eliminate reward hacking.
\newblock arXiv:2312.09244, 2023.

\bibitem[Diverse RM(2024)]{diverserm2024}
Multi-LoRA reward model ensembles with uncertainty weighting.
\newblock arXiv:2401.00243, 2024.

\bibitem[Efficient RM(2024)]{efficientrm2024}
Efficient reward model ensembles for robust RLHF.
\newblock arXiv:2401.16635, 2024.

% === Foundational / Representation ===
\bibitem[Holtzman et al.(2020)]{holtzman2020neural}
A.~Holtzman, J.~Buys, L.~Du, M.~Forbes, and Y.~Choi.
\newblock The curious case of neural text degeneration.
\newblock In \textit{ICLR}, 2020. arXiv:1904.09751.

\bibitem[Wang and Isola(2020)]{wang2020understanding}
T.~Wang and P.~Isola.
\newblock Understanding contrastive representation learning through alignment and uniformity on the hypersphere.
\newblock In \textit{ICML}, 2020.

\bibitem[Gao et al.(2022)]{gao2022contrastive}
T.~Gao, X.~Yao, and D.~Chen.
\newblock SimCSE: Simple contrastive learning of sentence embeddings.
\newblock In \textit{EMNLP}, 2022. arXiv:2104.08821.

\bibitem[Jing et al.(2019)]{jing2019understanding}
L.~Jing, S.~Vicol, and J.~Zhang.
\newblock Understanding dimensionality collapse in representation learning.
\newblock arXiv:1907.12009, 2019.

% === Scaling & Surveys ===
\bibitem[Scaling Laws Synthetic(2025)]{scaling2025synthetic}
Scaling laws for synthetic data training.
\newblock arXiv:2503.19551, 2025.

\bibitem[Synthetic CPT(2024)]{synthcpt2024}
Synthetic continued pretraining with selective accumulation.
\newblock arXiv:2409.07431, 2024.

\bibitem[Demystifying(2025)]{demystifying2025}
Demystifying synthetic data: A systematic study.
\newblock arXiv:2510.01631, 2025.

\bibitem[BeyondWeb(2025)]{beyondweb2025}
BeyondWeb: Trillion-scale synthetic pretraining experience.
\newblock arXiv:2508.10975, 2025.

\bibitem[QD Survey(2024)]{qd_survey2024}
Quality, diversity, and complexity: A survey of synthetic data evaluation.
\newblock arXiv:2412.02980, 2024.

\end{thebibliography}

%==========================================================
\appendix

\section{Complete Paper Inventory}
\label{sec:inventory}

Full list of 85 surveyed papers organized by category is available in the supplementary materials. See the panorama report at: \url{https://github.com/anonymous/collapse-survey-2026}.

\section{16-Point Prevention Checklist}
\label{sec:checklist}

\paragraph{Data Strategy}
\begin{enumerate}[label=C\arabic*]
    \item Use accumulation mode (retain real + append synthetic) --- \textit{Mandatory}
    \item Maintain real data fraction $\geq 0.618$ --- \textit{Strongly recommended}
    \item Use multi-source/multi-model synthetic generation --- \textit{Strongly recommended}
    \item Evaluate synthetic data diversity before training --- \textit{Recommended}
\end{enumerate}

\paragraph{Training Strategy}
\begin{enumerate}[label=C\arabic*, start=5]
    \item Apply truncated cross-entropy loss ($\gamma > 0.9$) --- \textit{Recommended}
    \item Consider negative guidance for diffusion models --- \textit{Recommended}
    \item Use entropy-based data selection --- \textit{Recommended}
    \item Apply regularization (KL/Wasserstein) for RLHF --- \textit{Recommended}
\end{enumerate}

\paragraph{Quality Control}
\begin{enumerate}[label=C\arabic*, start=9]
    \item Monitor synthetic data entropy (decrease = collapse signal) --- \textit{Mandatory}
    \item Track test loss trend (increase = collapse signal) --- \textit{Mandatory}
    \item Detect machine-generated text in training data --- \textit{Supplementary}
    \item Embed watermarks in synthetic data --- \textit{Supplementary}
\end{enumerate}

\paragraph{Warning Flags}
\begin{enumerate}[label=C\arabic*, start=13]
    \item Quality filtering must preserve diversity --- \textit{Warning}
    \item Synthetic data value is non-monotonic --- \textit{Warning}
    \item Verifiers have theoretical limits --- \textit{Warning}
    \item Even 1\% synthetic data can cause degradation --- \textit{Warning}
\end{enumerate}

\section{Statistical Analysis Details}
\label{sec:stats}

\paragraph{Golden Ratio Meta-Analysis.}
The 8 observed optimal mixing ratios are: $\{0.615, 0.620, 0.500, 0.60, 0.99, 0.65, 0.60, 0.55\}$. Testing against $H_0: \mu = 1/\varph = 0.618$:
\begin{itemize}[nosep]
    \item Sample mean: $\bar{\alpha}^* = 0.604$, SD $= 0.140$, SE $= 0.049$
    \item $t(7) = (0.604 - 0.618) / 0.049 = -0.286$, $p = 0.782$
    \item Excluding the outlier (classification, $\alpha^* = 0.99$): $\bar{\alpha}^* = 0.594$, SE $= 0.023$
    \item $t(6) = (0.594 - 0.618) / 0.023 = -1.043$, $p = 0.339$
    \item Cohen's $d = 0.93$ (large effect: mean is close to prediction)
    \item Bootstrap 95\% CI (10K resamples): [0.534, 0.676]
\end{itemize}
In all cases, we fail to reject $H_0$, supporting the golden ratio prediction.

\end{document}


